\subsection{Group Size, Consensus Quality, and Between-Team Variation}
\label{sec:group-size}

Early additions raise mean quality; later additions mainly improve the lower tail and reduce variation between teams. On the fixed panel of 171 Manual panoramas, each with 30--35 eligible annotators and 256 random entry orders, increasing team size from 3 to 6 raises mean $Q$ by 4.739 points (95\% path-bootstrap CI $[4.629,4.851]$). Increasing size from 20 to 30 adds only 0.500 points to the mean, but raises the within-image 10th percentile (P10) by 2.848 and reduces the within-image standard deviation by 2.031 (Figure~\ref{fig:group-size-171}, Table~\ref{tab:group-size-intervals}). P10 improves on 166 of 171 images, standard deviation declines on 168, and footprint variation declines on 170 (Supplement C5).

\begin{figure*}[t]
\centering
\setcounter{subfigure}{0}
\begin{minipage}{\textwidth}
\centering
{\footnotesize
\textcolor{figAll}{\rule[2pt]{10pt}{1.2pt}} All (171)\quad
\textcolor{paperBlue}{\rule[2pt]{10pt}{0.7pt}} Easy (116)\quad
\textcolor{figGreen}{\rule[2pt]{4pt}{0.7pt}\hspace{2pt}\rule[2pt]{4pt}{0.7pt}} Moderate (28)\quad
\textcolor{paperOrange}{\rule[2pt]{6pt}{0.7pt}\hspace{2pt}\rule[2pt]{1pt}{0.7pt}} Difficult (27)\par}
\vspace{3pt}
\ESPanel{.2414\linewidth}{panels/fig10/a.pdf}{Mean quality}%
\hfill%
\ESPanel{.2414\linewidth}{panels/fig10/b.pdf}{Lower-tail quality}%
\hfill%
\ESPanel{.2414\linewidth}{panels/fig10/c.pdf}{Between-team variation}%
\hfill%
\ESPanel{.2414\linewidth}{panels/fig10/d.pdf}{Complete outputs}

\end{minipage}
\caption{Group size affects mean Q, lower-tail Q, between-team variation and output completion differently. Curves use the same 171 images: 116 easy, 28 moderate and 27 difficult. Q summaries use scorable outputs; complete-and-evaluable rates use all 43,776 attempts per size. P10 is the image-specific 10th percentile, averaged over images. Shading marks 3--6 and 20--30 people. The at-least-half threshold admits exactly half of an even-sized team: the required support stays at two votes from three to four people, then rises to three at five people, contributing to odd--even alternation.}
\label{fig:group-size-171}
\end{figure*}


\begin{table}[t]
\centering\small
\caption{Changes from 20 to 30 people on the fixed 171-image panel, with equal image weights. Supplement D2 reports the quality levels.}
\label{tab:group-size-milestones}
\label{tab:group-size-intervals}
\setlength{\tabcolsep}{3pt}
\begin{tabular}{@{}lrrr@{}}
\toprule
Outcome & Change & Path interval & Building interval\\
\midrule
Mean Q & +0.500 & [0.439, 0.563] & [0.037, 1.042]\\
P10 & +2.848 & [2.620, 3.136] & [1.942, 3.970]\\
SD of Q & $-$2.031 & [$-$2.087, $-$1.962] & [$-$2.638, $-$1.503]\\
\bottomrule
\end{tabular}
\par\smallskip
\parbox{\columnwidth}{\footnotesize Intervals are approximate 95\% intervals. Path resampling measures Monte Carlo precision within the observed pools; building resampling retains the path estimates and resamples 22 buildings.}
\end{table}

Completion follows a separate trajectory. The complete-and-evaluable rate falls by 1.857 percentage points from 3 to 6 people, then recovers; from 20 to 30 it rises from 99.068\% to 99.488\% of 43,776 attempts per size. Quality summaries use scorable outputs, while completion uses every attempted team. Restricting the late-stage comparison to the 43,212 paths scorable at both sizes preserves the directions of mean, P10 and dispersion changes (Supplement D2).

Finite pools explain part of the narrowing variation: 30-person teams use 85.7\%--100\% of their candidate pools and increasingly share members. Yet removing the 12 images with exactly 30 candidates, where full-pool variation reaches zero, still gives a P10 gain of 2.942 points on the remaining 159 images, positive on 155. Capacity-stratified and full-pool results appear in Supplement C10.

Small aggregate changes can also conceal opposing image-level responses. Across 27 difficult panoramas, the 20-to-30-person mean change is $-0.174$ points, but 16 images improve and 11 deteriorate, with a mean absolute change of 4.772 points. The consensus continues to change as added observations alter support and median positions. A turn with four supporters is retained by an eight-person team but falls below the threshold when a ninth member omits it; new endpoint observations can also shift a retained turn without changing the corner count. Thus later additions stabilize most sampled teams while continuing to reshape contested layouts.

